% Rúbrica: Descripción de los datos (fuente, calidad y relevancia para el problema).
\section{Descripción de los datos}\label{sec:datos}

Los datos son reseñas públicas de Steam y la ficha de cada juego en la misma tienda. En cada
reseña, Steam guarda si el voto fue negativo y cuántos minutos llevaba jugados quien la escribió;
con esas dos cosas se define la señal de arrepentimiento temprano (sección~\ref{sec:objetivo}). El
modelo se entrena con \cifraJuegosEntrenamiento{} juegos y \cifraResenasEntrenamiento{} reseñas
(data-v1), y el sitio sirve un catálogo de \cifraJuegosCatalogo{} juegos y
\cifraResenasCatalogo{} reseñas (data-v3). No se usó Kaggle ni ningún conjunto de terceros.

\subsection{La API de Steam}\label{sec:datos-api}

La ingesta, el proceso que bajó los datos de Steam, usa dos servicios públicos de la
tienda\footnote{Una API (interfaz de programación) es una dirección web que devuelve datos en un
formato que un programa puede leer, en lugar de una página pensada para personas. Estos dos
servicios no piden clave.}. \texttt{appreviews} da, por reseña, el voto, los minutos jugados al
escribir y en total, el texto, las fechas y dos datos del autor: cuántos juegos tiene y cuántas
reseñas ha escrito. Se pidieron las más recientes, en inglés y de cualquier forma de compra, hasta
\cifraTopeIngesta{} por juego \parencite{valve_getreviews}. Hoy Valve marca ese servicio como
obsoleto y remite a \texttt{GetAppReviews} \parencite{valve_iuserreviews}, pero el proyecto no
depende de volver a bajarlo. \texttt{appdetails}, que Valve no documenta de forma oficial, da por
juego el precio y el descuento en pesos mexicanos, los géneros, la fecha de lanzamiento y la nota
de Metacritic\footnote{Metacritic promedia las calificaciones de la crítica especializada en una
nota de 0 a 100. Steam muestra esa nota en la página de algunos juegos.} tal como la muestra la
tienda: el proyecto nunca consulta Metacritic. En los \cifraJuegosExternos{} títulos externos, la
nota guardada coincide con la de la tienda en \cifraMetacriticVerificadoExternos{}\footnote{La
comparación está en \texttt{docs/evidencia/verificacion-40-steam.csv}.}.

\subsection{Releases y periodo}\label{sec:datos-releases}\label{sec:datos-periodo}

Cada corte de datos se publica como un release de GitHub con un tag fijo\footnote{Un release es
una versión publicada del repositorio, con archivos adjuntos. Su tag es el nombre que la
identifica, como \texttt{data-v1}, y no cambia: quien lo baje dentro de un año recibe los mismos
archivos.}, y antes de usar un archivo se compara su sha256 con el publicado\footnote{El sha256 es
una huella de 64 caracteres que se calcula sobre el contenido de un archivo; si cambia un solo
byte, cambia la huella. La comparación la hace \texttt{descargar\_verificado}, en
\texttt{backend/despliegue/utilidades.py}.}. Con data-v1 se toman todas las decisiones del modelo:
las variables, el entrenamiento y los umbrales. El release data-v2 suma \cifraJuegosExternos{}
juegos que el modelo nunca ve, y data-v3 es lo que sirve el sitio. Ningún release guarda el
identificador de Steam de quien escribió la reseña; el anexo~\ref{sec:anexo} da cada archivo con
su sha256.

Las reseñas de data-v1 se bajaron el \cifraDescargaHastaEntrenamiento{}, y las de los juegos que
agregó data-v2, hasta el \cifraDescargaHastaCatalogo{}, en UTC\footnote{UTC es el tiempo universal
coordinado. La Ciudad de México va seis horas atrás.}. La ingesta pide las más recientes hasta el
tope, y \cifraJuegosEnElTope{} de los \cifraJuegosEntrenamiento{} juegos lo alcanzan: un juego
popular lo llena en pocos días y uno de nicho necesita años, así que la muestra de cada juego
cubre de \cifraDiasCubiertosMin{} a \cifraDiasCubiertosMax{} días, con una mediana de
\cifraDiasCubiertosMediana{}. Por eso el \cifraDesdeJulioEntrenamientoPorResena{} de las reseñas
de entrenamiento es de julio de 2025 en adelante, con una cola que llega hasta el
\cifraResenaMasAntiguaEntrenamiento{} (detalle en el notebook 00, §3.1). La muestra retrata el
presente de cada juego: en uno viejo, la señal temprana viene de quien lo compra tarde. Es una
limitación de la fuente (sección~\ref{sec:limitaciones}) y no se corrige con pesos, porque no hay
forma de saber cómo era la mezcla de compradores en otra época.

\subsection{Calidad}\label{sec:datos-calidad}

Antes de explorar, el notebook 00 valida data-v1: tipos y nulos, llaves, rangos, textos repetidos
y la partición de validación (detalle en el notebook 00, §1). Las llaves y los rangos imposibles
no tienen un solo caso. Las dos columnas que definen la señal, el voto y los minutos al reseñar,
no tienen nulos, y todas las variables del juego que usa el modelo se conocen antes de comprar.
Tres hallazgos tocan al modelo, y ninguno se rellena en silencio. El precio de los
\cifraPagoSinPrecioEntrenamiento{} juegos de pago que llegaron sin precio se toma como 0, y la
ficha lo avisa. La falta de nota de los \cifraSinNotaEntrenamiento{} juegos sin Metacritic entra
al modelo como una variable propia. Y el 0 de los perfiles privados\footnote{En Steam, cada
persona decide si su perfil es público. Un perfil privado oculta la biblioteca, y la API lo
reporta con 0 juegos.}, que son el \cifraPrivadosEntrenamientoPorResena{} de las reseñas, se trata
como una bandera de privacidad y no como una biblioteca vacía. Los textos repetidos se marcan y se
conservan (sección~\ref{sec:procesamiento}).

El tope de \cifraTopeIngesta{} reseñas por juego es una decisión de diseño: así cada título pesa
casi lo mismo en el entrenamiento, y unos cuantos juegos muy populares no dominan la muestra.
Tiene otra consecuencia, los empates de tamaño entre juegos, que se trata con la validación
(sección~\ref{sec:modelacion-validacion}).

\subsection{Por qué estos datos sirven para la pregunta}\label{sec:datos-relevancia}

NexPlay quiere saber, antes de la compra, qué tan seguido un juego deja a sus compradores
arrepentidos en las primeras horas. Las reseñas las escriben personas que tienen el juego, y cada
una dice cuántos minutos llevaba jugados su autor: con eso se sabe si cae dentro o fuera de la
ventana de reembolso de dos horas. \texttt{appdetails} da lo que cualquiera ve en la tienda antes
de pagar, como el precio, el descuento y la nota de la crítica, y de ahí salen las variables del
modelo. También tienen límites (sección~\ref{sec:limitaciones}): las reseñas son en inglés y las
escriben compradores de todo el mundo, mientras que los precios son los de la tienda de México.
